\section{The original relation coefficients, their complete tail, and mixed source subtraction}
\label{sec:pgrt}

Retain the full quartet and every original map of PGS.1--41,
PGB.14--23 and PGMB.1--10. In particular
\[
\begin{gathered}
 h(s)=\prod_{r\in\{\rho,\bar\rho,1-\rho,1-\bar\rho\}}(s-r)^m,
 \quad \rho=\tfrac12+\delta+i\gamma,\quad0<\delta<\tfrac12,
 \quad\gamma>2,\quad m,k\geq1,\\
 v_h=2\xi/h,\quad\Phi'=h,\quad\Phi(0)=0,
 \quad\Phi(\tfrac12+it)=A+iq_\Phi(t),\\
 s_i=\tfrac12+it_i,\quad u=\sum_it_i,\quad S(u)=k/2+iu,
 \quad\tau_k=\sum_i\Phi(s_i),\quad w(t)=|v_h(\tfrac12+it)|^2/(2\pi),\\
 \chi(S)=\prod_{a,b=0}^k
 (S-k/2-(2a-k)\delta-i(2b-k)\gamma)^{1+k(m-1)},
 \quad q=[1+k(m-1)](k+1)^2,\quad E=\mathbb{C}[S]/(\chi).
\end{gathered}
\tag{PGRT.1}\label{eq:pgrt-original}
\]
The frame of $E$ is $1,S,\ldots,S^{q-1}$. The map $s:E\to\mathbb{C}[S]$
is the unique remainder of degree less than $q$ and $\iota P(u)=P(S(u))$.
The measure and full observation are
\[
\begin{gathered}
 r_k(u)=\int_{\sum t_i=u}(1+|\tau_k|^2)\prod_iw(t_i),
 \quad\mathscr L_k=L^2(r_k(u)\,du),\\
 F_b=\operatorname{rem}_{h(s_1),\ldots,h(s_k)}(\tau_k\Sigma^b),
 \quad \Sigma=\sum_i s_i,\quad
 z_b(u)=\int_{\sum t_i=u}\overline{\tau_k}F_b\prod_iw(t_i),\\
 g_x(u)=z(u)x/r_k(u),\quad F^0x=\iota sx-g_x.
\end{gathered}
\tag{PGRT.2}\label{eq:pgrt-f0}
\]
The fibre integral eliminates $t_k$ and uses $dt_1\cdots dt_{k-1}$;
at $k=1$ it is evaluation. Null exceptional sets in $z/r_k$ have
arbitrary values. PGS.39 proves $g:E\to\mathscr L_k$ is bounded.
All polynomial columns belong to $\mathscr L_k$ by the proved
exponential moments PGB.15. Inner products are conjugate linear
in their first variable. The original masses and full units remain
in every integral in this section.

\subsection{An exact analysis map for the unchanged relation source}

For each integer $j\geq q-1$ put $n=j+1-q$. There is a unique monic
polynomial $p_n$ of degree $n$ orthogonal to all lower polynomials
for the measure $|\chi(S(u))|^2r_k(u)\,du$. Indeed the Gram of
$1,S,\ldots,S^{n-1}$ is positive definite: a nonzero polynomial
cannot vanish almost everywhere for a density positive off a finite
set. Solving that Gram system for the lower coefficients of
$S^n$ gives existence and uniqueness, including $p_0=1$. Define
\[
 d_j=\chi p_{j+1-q},\quad
 a_j=\int|d_j(S(u))|^2r_k(u)\,du>0,\quad
 \alpha_jx=\langle\iota d_j,F^0x\rangle_{\mathscr L_k}.
\tag{PGRT.3}\label{eq:pgrt-relation-coefficients}
\]
These $d_j$ are exactly PGS.25: subtracting the old relation
projection from $\chi S^{j+1-q}$ leaves a monic relation of the
same degree, orthogonal to $\chi\mathcal P_{j-q}$; uniqueness
just proved identifies the two constructions. In particular
there is no change of the monic relation columns or their lengths.

The complete orthogonal family $\iota d_j$ spans a dense subspace
of $\mathscr L_k$. For specificity, PGB.15 proves an exponential
moment for $|\chi|^2r_k\,du$. Orthogonality to all its polynomials
forces the Fourier transform of the resulting finite complex
measure to have all derivatives zero on a nonempty analytic strip;
Fourier uniqueness, proved there by convolution with the exact
unit-mass Gaussian, forces that measure to vanish. Thus polynomials
are dense for this measure. Multiplication by $\chi(S(u))$ is
the onto isometry to $\mathscr L_k$, with inverse division off
its finite zero set, by the direct norm identity PGB.16. These
facts prove the asserted completeness of precisely these $d_j$.

Keep their norms in the sequence space
\[
 \ell^2(a^{-1})=
 \{(c_j)_{j\geq q-1}:\sum_{j\geq q-1}|c_j|^2/a_j<\infty\}.
\]
The exact analysis and synthesis maps are
\[
 \mathcal F:\mathscr L_k\longrightarrow\ell^2(a^{-1}),
 \quad f\longmapsto(\langle\iota d_j,f\rangle)_j,
 \qquad
 \mathcal F^{-1}c=\sum_{j=q-1}^{\infty}a_j^{-1}c_j\iota d_j.
\tag{PGRT.4}\label{eq:pgrt-analysis}
\]
Orthogonality proves that the squared norm of each finite sum is
$\sum|c_j|^2/a_j$, hence proves convergence for every displayed
sequence. Its coefficients are $c_j$, first for finite sequences
and then by continuity of each inner product. Conversely subtract
the synthesis of the coefficients of $f$; finite orthogonal
projection gives the Bessel bound, and the difference has zero
inner product with the complete family, hence is zero. This proves
both inverse identities, norm equality and zero kernels. No
orthogonal column has been rescaled in this construction.

Let $P_N^{\rm rel}$ denote orthogonal projection onto
$\iota(\chi\mathcal P_{N-q})$, with zero image for $N=q-1$.
The projection has the exact expression
\[
 P_N^{\rm rel}f=\sum_{j=q-1}^{N-1}a_j^{-1}
                 \langle\iota d_j,f\rangle\iota d_j,
 \qquad f_N=(I-P_N^{\rm rel})F^0.
\tag{PGRT.5}\label{eq:pgrt-projected-tail}
\]
To identify $f_N$ with PGMB.4, minimize the PGB.19 norm over
$P=sx+\chi Q$, $\deg P\leq N$. Its scalar component is
$F^0x+\iota\chi Q$, so its unique minimum is obtained by
subtracting $P_N^{\rm rel}F^0x$. The component has exactly the
form $\iota\widehat R_Nx-g_x$. This proves the displayed identity
for the original section rather than selecting another lift.

The complete residual form is therefore
\[
 T_N:=\widehat G_N-\Omega_k=f_N^*f_N
   =\sum_{j=N}^{\infty}Q_j,
 \qquad Q_j=a_j^{-1}\alpha_j^*\alpha_j,
 \qquad \Omega_k=Y_\perp^*Y_\perp>0.
\tag{PGRT.6}\label{eq:pgrt-complete-tail}
\]
For each $x$, PGRT.4--5 prove equality of the associated squared
norms. Polarization gives every matrix entry in the stated frame.
The series converges in operator norm: its positive tail has trace
equal to the sum of the squared tail norms of the $q$ fixed frame
vectors, which tends to zero by PGRT.4; its operator norm is at most
that trace. The exact kernel, without imposing additional analytic
injectivity, is
\[
 \ker f_N=\bigcap_{j\geq N}\ker\alpha_j
 =\{x\in E:F^0x\in\iota(\chi\mathcal P_{N-q})\}
 =\{x:g_x=\iota P\text{ for some }\deg P\leq N,
                                    \ [P]=x\}.
\tag{PGRT.7}\label{eq:pgrt-tail-kernel}
\]
The first equality is the positive norm sum. The second is the
projection kernel, and the third follows by subtracting the stated
relation from $sx$. This is the full supported-zero fibre of the
tail map, while the analysis isometry itself has zero kernel.

Every row in this expansion has its original arithmetic integral:
\[
 (\alpha_j)_b=\int_{\mathbb{R}}\overline{d_j(S(u))}
                  \bigl(S(u)^b r_k(u)-z_b(u)\bigr)\,du.
\tag{PGRT.8}\label{eq:pgrt-explicit-row}
\]
It is absolutely convergent by Cauchy--Schwarz. Expanding the
polynomials and inserting the PGS.35 and PGS.38 convolution
formulas expresses every term as a finite sum of products of
$M^{(h)}_{ab}=\int\bar s^as^b|v_h(s)|^2dt/(2\pi)$.
The coefficients of $d_j$ are obtained by the positive Gram
system in PGRT.3, whose entries are those same finite moments.
Since $\widehat R_j-s$ is an old relation, its scalar component
pairs to zero with $d_j$. Thus this row is exactly
$\beta_j-d_j^*Z_{j+1}$ in PGS.25, including its negative cross
term. This proves an explicit finite arithmetic route to each
coefficient, even though the complete tail is infinite.

\subsection{A strict mixed subtraction in the endpoint upper bound}

Retain the original Hilbert maps
$D_{\rm mix}=\mathscr D_k^{\rm mix}:E\to\mathscr K_k^{\rm mix}$
and $L_{\rm low}:E\to\mathscr H_k$ of PGMB.2 and PGMB.6.
For a fixed integer $M\geq2q$ define
\[
 H_M=\widehat G_M=K_{\rm low}+Z_M^{\rm mix\,*}Z_M^{\rm mix},
 \qquad Z_M^{\rm mix}x=(D_{\rm mix}x,f_Mx).
\tag{PGRT.9}\label{eq:pgrt-mixed-pair}
\]
Here $Z_M^{\rm mix}$ maps into the ordered Hilbert direct sum
$\mathscr K_k^{\rm mix}\oplus\mathscr L_k$; it is not the
finite cross-moment matrix $Z_M$ of PGS.5. The identity is
exactly PGMB.5. The operator
$I+Z_M^{\rm mix}K_{\rm low}^{-1}Z_M^{\rm mix\,*}$ is invertible:
on the orthogonal complement of the finite-dimensional range
of $Z_M^{\rm mix}$ it is the identity, and on that range it is
positive with all eigenvalues at least one. Direct multiplication
of both sides by $H_M$ proves
\[
 H_M^{-1}=K_{\rm low}^{-1}-K_{\rm low}^{-1}Z_M^{\rm mix\,*}
 (I+Z_M^{\rm mix}K_{\rm low}^{-1}Z_M^{\rm mix\,*})^{-1}
                     Z_M^{\rm mix}K_{\rm low}^{-1}.
\tag{PGRT.10}\label{eq:pgrt-inverse-subtraction}
\]
For the multiplication, write $K=K_{\rm low}$ and $Z=Z_M^{\rm mix}$;
the leftover terms factor as
$Z^*ZK^{-1}-Z^*(I+ZK^{-1}Z^*)(I+ZK^{-1}Z^*)^{-1}ZK^{-1}=0$.
Thus all ordered products in the inverse formula are retained.

For $q-1\leq j\leq2q-1$ its scalar consequence is the exact
nonnegative value
\[
\begin{aligned}
 u_{j,M}&=a_j^{-1}\alpha_jH_M^{-1}\alpha_j^*\\
 &=a_j^{-1}\alpha_jK_{\rm low}^{-1}\alpha_j^*
  -a_j^{-1}\left\|
 (I+Z_M^{\rm mix}K_{\rm low}^{-1}Z_M^{\rm mix\,*})^{-1/2}
                   Z_M^{\rm mix}K_{\rm low}^{-1}\alpha_j^*
                         \right\|^2.
\end{aligned}
\tag{PGRT.11}\label{eq:pgrt-strict-subtraction}
\]
If $k\geq2$ and $\alpha_j\ne0$, the subtracted term is strictly
positive. Indeed PGM.6--9 proves $D_{\rm mix}$ injective, so
$Z_M^{\rm mix}K_{\rm low}^{-1}\alpha_j^*\ne0$, and the displayed
positive inverse square root has zero kernel. When $k=1$ the mixed
component is zero, but its finite tail component is still present;
no strictness is asserted from the vanishing component. If
$\alpha_j=0$, both terms are exactly zero. Consequently the strict
mixed contribution gives a calculated subtraction at each nonzero
relation row, not merely a positive addition to an unrelated norm.

Let $\widehat{\mathcal B}_k$ denote exactly the four signs of PGS.29,
and retain $\omega_{q-1}=\omega_{2q-1}=1$ and $\omega_j=2$
for $q\leq j\leq2q-2$. Set
\[
 l_j=a_j^{-1}\alpha_j\widehat G_{q-1}^{-1}\alpha_j^*,
 \quad\eta_j=a_j^{-1}\alpha_j\widehat G_{j+1}^{-1}\alpha_j^*.
\]
PGRT.6 gives $\widehat G_j=\widehat G_{j+1}+Q_j$.
Congruence with $\widehat G_{j+1}^{-1/2}$ gives a rank-one
addition with eigenvalues $1+\eta_j,1,\ldots,1$; extending the
nonzero update vector to an orthonormal basis proves this statement,
and the zero vector gives the identity matrix. Therefore
\[
\begin{gathered}
 \log\frac{\det\widehat G_j}{\det\widehat G_{j+1}}
       =\log(1+\eta_j),\quad
 \widehat{\mathcal B}_k=
                  \sum_{j=q-1}^{2q-1}\omega_j\log(1+\eta_j),\\
 L_k^{\rm rel}:=\sum_j\omega_j\log(1+l_j)
 \ \leq\ \widehat{\mathcal B}_k\ \leq\
 U_{k,M}^{\rm rel}:=\sum_j\omega_j\log(1+u_{j,M}).
\end{gathered}
\tag{PGRT.12}\label{eq:pgrt-endpoint-bounds}
\]
For the inequalities use
$H_M\preceq\widehat G_{j+1}\preceq\widehat G_{q-1}$.
For positive forms inverse order reverses: conjugate by the square
root of the smaller form, diagonalize the resulting matrix with
eigenvalues at least one, invert them, and undo the congruence.
Thus $l_j\leq\eta_j\leq u_{j,M}$, and increasing logarithms
prove PGRT.12. Telescoping along the two exact endpoint intervals
proves the weights. For $q=1$ the interior is empty and the two
steps remain distinct, so no repeated endpoint is dropped before
the cancellation. The actual quartet has its full $q$ in PGRT.1.

The choice $M=2q$ gives the strongest upper bound among these
common later-window comparisons, because $H_M$ decreases as $M$
increases. All $H_M$ in PGRT.11--12 are finite arithmetic Gram
matrices computed by PGS.9--13. At $M\to\infty$ they tend to
$\Omega_k$, and $f_M\to0$, so PGRT.10--11 also prove the exact
limiting subtraction with $Z_\infty^{\rm mix}=D_{\rm mix}$ and
$H_\infty=\Omega_k$. No rate in the separate variable $k$ has
been presumed. The strongest available local upper bound can
retain the minimum of PGRT.12 and the proved PGMB.10 value
$t_{q-1}+t_q$, where $t_N=\log\det(I+\Omega_k^{-1/2}T_N\Omega_k^{-1/2})$.

\subsection{Return to the original arithmetic endpoint}

Write $\Gamma_k^{\rm gr}=\widehat{\mathcal B}_k-\mathcal B_k$
for the literal signed augmentation, keeping the previously proved
finite original-source enclosure $g_s^-\leq\Gamma_k^{\rm gr}\leq g_s^+$
of ACM27. This enclosure is the output of its full signed projector
calculation and its proved finite Neumann remainder, not an assumed
estimate. Subtracting those two inequalities from PGRT.12 yields
\[
 L_k^{\rm rel}-g_s^+\leq\mathcal B_k
                  \leq U_{k,2q}^{\rm rel}-g_s^-.
\tag{PGRT.13}\label{eq:pgrt-arithmetic-return}
\]
Intersecting this interval with the already proved Gamma-to-arithmetic
interval $[\beta_{p,\Gamma}^-,\beta_{p,\Gamma}^+]$ gives
\[
\begin{aligned}
 b_{p,s}^{\rm rel,-}
 &=\max\{\beta_{p,\Gamma}^-,L_k^{\rm rel}-g_s^+\},\\
 b_{p,s}^{\rm rel,+}
 &=\min\{\beta_{p,\Gamma}^+,U_{k,2q}^{\rm rel}-g_s^-\},\qquad
 b_{p,s}^{\rm rel,-}\leq\mathcal B_k\leq b_{p,s}^{\rm rel,+}.
\end{aligned}
\tag{PGRT.14}\label{eq:pgrt-current-interval}
\]
Each input is a defined original arithmetic matrix or a bound already
proved for it. The interval is nonempty because it contains the same
$\mathcal B_k$. At HC9--10 use the unchanged $T_k=4q\log(D_hk)$:
subtract $T_k$ from both endpoints for $\mathcal R_k$, and subtract
$\mathcal B_k^\Gamma$ for $\Delta_k^\Gamma$. Intersect with the
earlier lower and upper inequalities before these subtractions.
Their validity follows directly from equality of the retained scalar
and does not discard the negative $g_s^-$ term. The phase propagation
HC14g then uses the resulting arithmetic upper endpoint in each of
its three increasing scalar functions. This is the exact return map
of the new relation calculation into the existing endpoint program.

All analysis, projection, synthesis and quotient maps in this section
act coordinatewise on every fixed original support mask. Their
represented-zero fibres are products of the displayed linear kernels.
A present empty tuple maps to a present empty tuple, and external
absence maps to external absence with singleton inverse image.
None of the Hilbert completions identifies those two fibres.
The unresolved growing-$k$ analytic calculation is now the actual
size of the rows PGRT.8 in the full inverses PGRT.11--12 together
with the signed arithmetic return PGRT.13; the complete identities
here neither remove that calculation nor substitute an assumed bound.
